% Vietnamese translation for OI Public Library
% Source-PDF: IOI中国国家候选队论文/2000/李刚论文.pdf
\documentclass[11pt]{article}
\usepackage{oipl-vn}

\title{Bàn sâu về quy hoạch động}
\author{Li Gang}
\date{}

\begin{document}
\maketitle

\origtitle{An in-depth discussion of dynamic programming}
\sourcepath{IOI China candidate team papers / 2000 / Li Gang.pdf}

\paragraph{Từ khóa.}
Quy hoạch động; trạng thái.

\section*{Tóm Tắt}

Bài viết này thảo luận một kỹ thuật giải bài rất hiệu quả: quy hoạch động.
Hiệu suất giải bài cao của kỹ thuật này luôn được quan tâm nhiều. Trước hết,
bài viết bàn về nền tảng lý thuyết của quy hoạch động, và nêu hai tiền đề để
một bài toán có thể giải bằng quy hoạch động: cấu trúc con tối ưu và tính
không hậu hiệu. Tiếp đó, bài viết xét hai vấn đề thường gặp trong ứng dụng
thực tế: chọn trạng thái và lưu trữ trạng thái trong quy hoạch động. Từ các
thảo luận ấy, bài viết dẫn đến phương pháp tư duy cơ bản của quy hoạch động:
``không làm lại việc đã làm'', cũng như nguyên nhân khiến kỹ thuật quy hoạch
động có tốc độ đáng kinh ngạc khi giải bài: ``xử lý phần dư thừa trong việc
xét các khả năng, đạt tới giới hạn về tốc độ''. Cuối cùng, bài viết trình bày
các bước tổng quát để giải bài quy hoạch động: suy nghĩ, lập kế hoạch, và
thực hiện.

\tableofcontents

\section{Dẫn nhập}

Trong các kỳ thi tin học, đặc biệt trong vài năm gần đây, quy hoạch động
thường xuyên được nhắc đến như một công cụ giải bài. Phạm vi ứng dụng của nó
ngày càng rộng, mức độ ứng dụng cũng ngày càng sâu. Vậy quy hoạch động khác
những thuật toán khác ở điểm nào? Giá trị ứng dụng cụ thể của nó là gì? Bài
viết này sẽ thảo luận các câu hỏi đó.

Trước hết, ta nhận biết quy hoạch động qua một bài toán cụ thể.

\paragraph{Ví dụ 1.}
Hình 1 cho một bản đồ. Mỗi đỉnh trên bản đồ biểu diễn một thành phố, đường nối
giữa hai thành phố biểu diễn đường đi, và số trên đường nối biểu diễn chiều
dài đường. Nay ta muốn đi từ thành phố \(A\) đến thành phố \(E\). Đi thế nào
để quãng đường ngắn nhất, và độ dài ngắn nhất là bao nhiêu?

Giả sử:
\[
\begin{array}{ll}
\mathrm{Dis}[X] & \text{là độ dài đường ngắn nhất từ thành phố } X
                  \text{ đến } E,\\
\mathrm{Map}[I,J] & \text{là khoảng cách giữa hai thành phố } I,J.
\end{array}
\]
Nếu \(\mathrm{Map}[I,J]=0\) thì hai thành phố không thông nhau.

Bài toán này có thể làm bằng tìm kiếm; chương trình rất dễ viết:

\begin{verbatim}
Var
  Se: tap thanh pho chua tham;

Function Long(Who: thanh pho hien tai): Integer;
Begin
  If Who=E Then Search:=0
  Else
  Begin
    Min:=Maxint;
    For I qua moi thanh pho Do
      If (Map[Who,I]>0) And (I In Se) Then
      Begin
        Se:=Se-[I];
        J:=Map[Who,I]+Long(I);
        Se:=Se+[I];
        If J<Min Then Min:=J;
      End;
    Long:=Min;
  End;
End;

Begin
  Se:=tap moi thanh pho tru A;
  Dis[A]:=Long(A);
End.
\end{verbatim}

Hiệu suất của chương trình này ra sao? Ta thấy mỗi lần, trừ các thành phố đã
thăm, mọi thành phố khác đều phải được xét, nên độ phức tạp thời gian là
\(O(N!)\). Đây là một thuật toán cấp mũ. Vậy có thuật toán tốt hơn không?

Trước tiên hãy quan sát thuật toán này. Khi tìm đường ngắn nhất từ \(B_1\) đến
\(E\), ta đã tìm đường ngắn nhất từ \(C_2\) đến \(E\); khi tìm đường ngắn nhất
từ \(B_2\) đến \(E\), ta lại tìm một lần nữa đường ngắn nhất từ \(C_2\) đến
\(E\). Nghĩa là đường ngắn nhất từ \(C_2\) đến \(E\) bị tính hai lần. Tương tự,
trong quá trình tìm đường ngắn nhất từ \(C_1,C_2\) đến \(E\), đường ngắn nhất
từ \(D_1\) đến \(E\) cũng bị tính hai lần. Trong toàn bộ chương trình, đường
ngắn nhất từ \(D_1\) đến \(E\) bị tính bốn lần. Lãng phí lớn biết bao! Nếu
trong quá trình giải, ta đồng thời ghi lại độ dài đường ngắn nhất đã tìm được
và gọi lại bất cứ lúc nào thì sẽ tiện lợi hơn nhiều.

Từ đó sinh ra một ý tưởng mới: lần lượt suy ra từng giá trị \(\mathrm{Dis}\)
từ sau ra trước, cho đến khi suy ra \(\mathrm{Dis}[A]\). Ý tưởng này quả thật
tốt, nhưng khoan đã: thế nào là ``từ sau ra trước''?

``Sau'' và ``trước'' là thứ tự do ta tự gán cho các thành phố. Khi thứ tự của
hai thành phố \(I,J\) được định là \(I\) ``trước'', \(J\) ``sau'', thì phải
thỏa điều kiện:
\[
  I,J \text{ không thông nhau, hoặc }
  \mathrm{Dis}[I]+\mathrm{Map}[I,J]\ge \mathrm{Dis}[J].
\]
Nếu \(I,J\) thông nhau và
\(\mathrm{Dis}[I]+\mathrm{Map}[I,J]<\mathrm{Dis}[J]\), điều đó cho thấy
\(\mathrm{Dis}[J]\) còn có phương án tốt hơn; nhưng \(J\) lại nằm sau \(I\), nên
phương án ấy không thể được suy ra, làm ảnh hưởng đến lời giải cuối cùng. Vậy
ta chia thứ tự trước sau như thế nào?

Như Hình 2, các thành phố được chia thành nhiều giai đoạn. Ta có thể dùng
giai đoạn để biểu diễn thứ tự của mỗi thành phố, vì cách chia giai đoạn có
hai tính chất sau:
\begin{enumerate}
  \item Giá trị ở giai đoạn \(I\) chỉ liên quan đến giai đoạn \(I+1\), và giá
  trị ở giai đoạn \(I\) chỉ gây ảnh hưởng đến giá trị ở giai đoạn \(I-1\).
  \item Thứ tự giữa các giai đoạn là xác định; không thể đổi chỗ tùy ý hai
  giai đoạn bất kỳ.
\end{enumerate}
Từ hai tính chất này có thể suy ra rằng dùng giai đoạn làm thứ tự ``trước'' và
``sau'' thỏa điều kiện vừa nêu. Vì vậy, ta dùng giai đoạn làm thứ tự của mỗi
thành phố, rồi truy ngược từ giai đoạn 3 đến giai đoạn 1, sau đó suy ra
\(\mathrm{Dis}[A]\).

Công thức là
\[
  \mathrm{Dis}[X]=
  \min\{\mathrm{Dis}[Y]+\mathrm{Map}[X,Y]:
  Y \text{ là thành phố ở giai đoạn kế tiếp và thông với } X\}.
\]
Có thể xem \(E\) là giai đoạn 4 và \(A\) là giai đoạn 0.

Chương trình:
\begin{verbatim}
Dis[E]=0
For X= moi thanh pho tu giai doan 3 xuong giai doan 0 Do
Begin
  Dis[X]:=Maxint;
  For Y= moi thanh pho o giai doan ke tiep cua X Do
    If Dis[Y]+Map[X,Y]<Dis[X] Then
      Dis[X]:=Dis[Y]+Map[X,Y];
End.
\end{verbatim}

Độ phức tạp thời gian của chương trình này là \(O(N^2)\), nhỏ hơn rất nhiều
so với \(O(N!)\) của chương trình trước. Thuật toán thứ hai chính là thuật
toán quy hoạch động.

\section{Nền tảng lý thuyết của quy hoạch động}

Qua ví dụ trên, ta đã có nhận thức ban đầu về quy hoạch động: nó xử lý một
``bài toán quyết định nhiều giai đoạn''. Bây giờ ta định nghĩa cụ thể một số
khái niệm.

\term{Trạng thái} (state): biểu diễn tính chất của sự vật, là đại lượng mô tả
``đơn vị'' trong quy hoạch động. Chẳng hạn trong Ví dụ 1, các thành phố
\(A,B_1,D_2,E\) đều là những trạng thái riêng lẻ.

\term{Giai đoạn} (stage): một tập các trạng thái có tính chất gần nhau và có
thể xử lý đồng thời. Thông thường, một bài toán có thể được chia thành vài
giai đoạn theo thứ tự xử lý. Ví dụ 1 gồm ba giai đoạn, trong đó giai đoạn 1
chứa hai trạng thái \(B_1,B_2\). Thực ra giai đoạn chỉ đánh dấu những trạng
thái có cùng cách xử lý và không phụ thuộc vào thứ tự xử lý lẫn nhau. Một
giai đoạn có thể chứa nhiều trạng thái, cũng có thể chỉ chứa một trạng thái;
quan hệ giữa chúng khá giống quan hệ giữa phân tử và nguyên tử.

\term{Phương trình chuyển trạng thái}: quy luật biến đổi để trạng thái của
giai đoạn trước chuyển thành trạng thái của giai đoạn sau. Đây là phương trình
liên quan đến trạng thái của hai giai đoạn kề nhau, và là trung tâm của quy
hoạch động.

\term{Quyết định} (decision): một hành động lựa chọn nào đó được đưa ra ở mỗi
giai đoạn. Đó chính là lựa chọn mà chương trình của ta cần thực hiện.

Bài toán do quy hoạch động xử lý là một bài toán quyết định nhiều giai đoạn:
thường bắt đầu từ trạng thái đầu, thông qua việc chọn quyết định ở các giai
đoạn giữa, rồi đi đến trạng thái kết thúc. Các quyết định ấy tạo thành một
dãy quyết định, đồng thời xác định một tuyến hoạt động để hoàn thành toàn bộ
quá trình, thường là tuyến hoạt động tối ưu. Hình 3 minh họa quan hệ này.

Nhìn tổng thể, một thuật toán quy hoạch động nên chứa các quan hệ và thuộc
tính cơ bản vừa nêu.

Vậy bài toán quyết định như thế nào mới có thể chia giai đoạn và giải bằng
quy hoạch động? Ta nói rằng chỉ khi một bài toán quyết định có các tính chất
sau thì mới nên cân nhắc chia giai đoạn và dùng quy hoạch động.

\subsection{Cấu trúc con tối ưu}

Nghĩa là lời giải tối ưu của một bài toán chỉ phụ thuộc vào lời giải tối ưu
của các bài toán con; các lời giải không tối ưu không ảnh hưởng đến việc giải
bài toán. Trong Ví dụ 1, nếu muốn tìm giá trị nhỏ nhất của \(\mathrm{Dis}[B_1]\),
ta chỉ cần biết giá trị nhỏ nhất của \(C_1,C_2,C_3\), không cần xét các đường
khác, vì các đường không tối ưu chắc chắn không ảnh hưởng đến kết quả của
\(\mathrm{Dis}[B_1]\). Xét thêm một bài toán:

\paragraph{Ví dụ 2.}
Có bốn điểm \(A,B,C,D\), như Hình 4. Giữa hai điểm kề nhau có hai đường đi
\(C_{2k},C_{2k-1}\) \((1\le k\le 3)\). Số phía trên đường nối biểu diễn độ
dài đường. Ta định nghĩa trong mọi đường từ \(A\) đến \(D\), đường có số dư
nhỏ nhất khi lấy độ dài chia cho 4 là đường tối ưu. Hãy tìm một đường tối ưu.

Với bài này, nếu vẫn giải theo cách vừa dùng thì sẽ sai. Chẳng hạn, theo tư
duy của Ví dụ 1, giá trị tối ưu của \(A\) có thể xác định từ giá trị tối ưu
của \(B\); giá trị tối ưu của \(B\) là 0, nên giá trị tối ưu của \(A\) sẽ là
2. Nhưng thực tế, đường \(C_1-C_3-C_5\) cho giá trị tối ưu là 0. Vì vậy,
đường tối ưu của \(B\) không phải đường con của đường tối ưu của \(A\). Nói
cách khác, giá trị tối ưu của \(A\) không do giá trị tối ưu của \(B\) quyết
định, nên bài toán không có cấu trúc con tối ưu.

Từ đó có thể thấy không phải mọi bài toán quyết định đều có thể giải bằng quy
hoạch động. Chỉ khi một bài toán thể hiện cấu trúc con tối ưu thì quy hoạch
động mới có thể là một phương pháp thích hợp.

\subsection{Tính không hậu hiệu}

Sau khi một bài toán được chia giai đoạn, trạng thái trong giai đoạn \(I\) chỉ
có thể thu được từ trạng thái trong giai đoạn \(I+1\) thông qua phương trình
chuyển trạng thái, không liên quan đến các trạng thái khác, đặc biệt không
liên quan đến các trạng thái chưa xảy ra. Đây là tính không hậu hiệu. Thực ra,
nếu ta định nghĩa trạng thái của bài toán là các đỉnh trong đồ thị, chuyển
giữa hai trạng thái là cạnh, và lượng tăng của trọng số trong quá trình
chuyển là trọng số cạnh, thì bài toán thực chất là tìm đường giữa hai đỉnh
trong một đồ thị có hướng, không chu trình, có trọng số. Vì có tính không hậu
hiệu nên không có chu trình; nếu có chu trình thì dù chia giai đoạn thế nào
cũng có thể xuất hiện hậu hiệu. Nói cách khác, đồ thị ấy có thể sắp thứ tự tô
pô; ít nhất ta có thể chia giai đoạn theo thứ tự tô pô của nó.

Xét một ví dụ cụ thể.

\paragraph{Ví dụ 3.}
Bài toán người bán hàng Euclid yêu cầu, với \(n\) điểm cho trước trên mặt
phẳng, xác định một hành trình khép kín đi qua các điểm ấy. Hình 5(a) cho lời
giải của bài toán bảy điểm. Bài toán hành trình bitonic là một dạng đơn giản
hơn của bài toán người bán hàng Euclid: hành trình xuất phát từ điểm ngoài
cùng bên trái, đi nghiêm ngặt từ trái sang phải đến điểm ngoài cùng bên phải,
rồi lại đi nghiêm ngặt từ phải sang trái về điểm xuất phát; cần tìm độ dài
đường đi ngắn nhất. Hình 5(b) cho lời giải của bài toán bảy điểm.

Hai bài toán này nhìn có vẻ rất giống nhau, nhưng về bản chất lại khác nhau.
Để tiện thảo luận, ta đánh số các đỉnh. Vì hành trình chắc chắn đi qua đỉnh
ngoài cùng bên phải, là đỉnh 7, nên một đường \((P_1-P_2)\) có thể xem là sự
kết hợp của hai đường \((P_1-7)\) và \((P_2-7)\). Do đó trạng thái của bài
này có thể biểu diễn bằng dạng kết hợp của hai đường. Ta có thể xếp các trạng
thái có cùng hai đỉnh xuất phát của hai đường vào một giai đoạn, ký hiệu là
giai đoạn \([P_1,P_2]\).

Đối với bài toán hành trình bitonic, nếu giai đoạn \([P_1,P_2]\) có thể được
suy ra từ giai đoạn \([Q_1,Q_2]\), thì điều kiện bắt buộc là
\(P_1<Q_1\) hoặc \(P_2<Q_2\). Ví dụ, đường trong giai đoạn \([3,4]\) có thể
thu được từ đường trong giai đoạn \([3,5]\) bằng cách thêm cạnh \(4-5\);
nhưng trạng thái của giai đoạn \([3,5]\) lại không thể thu được từ trạng thái
của giai đoạn \([3,4]\), vì hành trình bitonic yêu cầu phải đi nghiêm ngặt từ
trái sang phải. Do đó nếu ta đã biết trạng thái trong giai đoạn \([3,4]\),
thì trạng thái trong giai đoạn \([3,5]\) tất yếu đã biết. Vì vậy, ta nói bài
toán bitonic thỏa nguyên tắc không hậu hiệu và có thể giải bằng quy hoạch
động; chương trình xem Pro\_3\_1.Pas.

Đối với bài toán người bán hàng Euclid, giữa các giai đoạn không có một thứ
tự tất yếu nào. Chẳng hạn đường
\(\{3-2-5-7,\ 4-6-7\}\) thuộc giai đoạn \([3,4]\), có thể suy ra từ đường
\(\{2-5-7,\ 4-6-7\}\) thuộc giai đoạn \([2,4]\); còn đường
\(\{2-3-6-7,\ 4-5-7\}\) thuộc giai đoạn \([2,4]\), có thể suy ra từ đường
\(\{3-6-7,\ 4-5-7\}\) thuộc giai đoạn \([3,4]\). Nếu dùng đỉnh để biểu diễn
giai đoạn và dùng quan hệ suy ra để biểu diễn cạnh, thì quan hệ tương ứng
giữa hai giai đoạn \([3,4]\) và \([2,4]\) giống Hình 6. Ta thấy rất rõ rằng
quan hệ giữa hai giai đoạn này có hậu hiệu, vì trong đồ thị có chu trình. Với
bài toán này, ta không thể giải như bài trước. Thực tế, đây là một bài toán
NP-đầy đủ, thời gian giải rất có khả năng là cấp mũ.

Vì vậy, một điều kiện phán đoán rất quan trọng để xem một bài toán có thể
giải bằng quy hoạch động hay không là: nó có hậu hiệu hay không. Khi phán
đoán điều này, một cách hiệu quả là xem các giai đoạn của bài toán như đỉnh,
quan hệ giữa các giai đoạn như cạnh có hướng, rồi kiểm tra đồ thị có hướng ấy
có phải là đồ thị không chu trình hay không, tức có thể sắp thứ tự tô pô hay
không.

Từ các phân tích trên, ta có thể rút ra một số phương pháp và bước cơ bản để
giải bài quy hoạch động:
\begin{enumerate}
  \item Xác định đối tượng nghiên cứu của bài toán, tức xác định trạng thái.
  \item Chia giai đoạn, xác định phương trình chuyển trạng thái giữa các giai
  đoạn.
  \item Xem bài toán hiện tại có thể giải bằng quy hoạch động hay không:
  kiểm tra bài toán có cấu trúc con tối ưu hay không, và có tính không hậu
  hiệu hay không.
  \item Nếu phát hiện bài toán hiện tại chưa thể giải bằng quy hoạch động,
  cần điều chỉnh các định nghĩa và cách chia tương ứng để đạt đến dạng có thể
  giải bằng quy hoạch động.
\end{enumerate}

Trên đây chỉ là quá trình tư duy quy hoạch động trong trường hợp thông
thường. Một số bài đơn giản có thể thao tác tuần tự theo các bước ấy, nhưng
đa số bài quy hoạch động, đặc biệt là bài trong thi tin học, kiểm tra nhiều
mảng kiến thức và có kỹ thuật vận dụng linh hoạt. Dưới đây ta thảo luận một
số điểm trọng yếu và khó trong ứng dụng quy hoạch động.

\section{Ứng dụng thực tế của quy hoạch động}

Khi đánh giá một thuật toán, tiêu chuẩn chủ yếu không ngoài hai chỉ tiêu: thời
gian và không gian. Thời gian của đa số thuật toán quy hoạch động là cấp đa
thức; so với thời gian cấp mũ của thuật toán tìm kiếm giải cùng bài toán,
thời gian của quy hoạch động là rất ít. Vì vậy trong ứng dụng thực tế, ta
hiếm khi cân nhắc vấn đề thời gian của quy hoạch động; vấn đề thường gặp nhất
là không gian, tức chọn và lưu trữ trạng thái.

\subsection{Chọn trạng thái}

Đối với một thuật toán quy hoạch động, việc chia giai đoạn không quá quan
trọng, vì giai đoạn chỉ là tập các trạng thái có thể xử lý giống nhau; ta hoàn
toàn có thể định nghĩa mỗi trạng thái riêng lẻ là một giai đoạn. Vì vậy, việc
chọn trạng thái đóng vai trò quyết định đối với toàn bộ cách xử lý bài toán.
Trạng thái được chọn phải thỏa hai điểm sau:
\begin{enumerate}
  \item Trạng thái phải mô tả đầy đủ tính chất của sự vật; hai sự vật khác
  nhau phải có trạng thái khác nhau.
  \item Phải tồn tại phương trình chuyển giữa các trạng thái, để ta có thể
  chuyển dần từ trạng thái ứng với bài toán ban đầu đến trạng thái ứng với
  bài toán kết thúc.
\end{enumerate}

Trạng thái là đại lượng mô tả tính chất của sự vật, vì vậy ta cần dựa vào
tiêu chuẩn này và phân tích cụ thể theo yêu cầu của đề bài. Xét ví dụ sau.

\paragraph{Ví dụ 4.}
Có một công ty vận chuyển bò, cần vận chuyển bò giữa bảy nông trại
\(A,B,C,D,E,F,G\). Xe xuất phát từ \(A\), sau khi hoàn thành nhiệm vụ phải
quay về \(A\). Mỗi lần xe chỉ chở được một con bò. Mỗi nhiệm vụ vận chuyển
được cho bởi một cặp chữ cái, lần lượt biểu thị bò được chuyển từ nông trại
nào đến nông trại nào. Tổng số nhiệm vụ là \(N\), \(1\le N\le 12\). Biết
khoảng cách giữa các nông trại, hãy tìm một lộ trình ngắn nhất để hoàn thành
mọi nhiệm vụ.

Ta phân tích cách giải. Mục tiêu của ta là hoàn thành một số nhiệm vụ. Vì xe
mỗi lần chỉ chở được một con bò, ta chỉ có thể xử lý các nhiệm vụ lần lượt.
Do đề bài yêu cầu lộ trình ngắn nhất để hoàn thành các nhiệm vụ, nên trong
quá trình hoàn thành nhiệm vụ, mọi đoạn đường ta đi đều là đường ngắn nhất.
Nói cách khác, một khi đã xác định thứ tự hoàn thành nhiệm vụ
\(P_1-P_2-\cdots-P_N\), thì thời gian ngắn nhất của lộ trình ấy cũng xác định.
Mục tiêu của ta là xác định thứ tự hoàn thành \(N\) nhiệm vụ. Vì vậy, đối
tượng nghiên cứu của bài toán là một tập nhiệm vụ \(S\).

Vậy trạng thái nên định nghĩa thế nào? Vì đối tượng nghiên cứu là tập hợp, ta
tự nhiên nghĩ đến cách mô tả quan hệ phần tử trong tập. Có thể định nghĩa
trạng thái là một mảng \((t_1,t_2,\ldots,t_N)\), trong đó mỗi \(t_i\) bằng 0
nếu nhiệm vụ \(i\) không thuộc tập hiện tại, và bằng 1 nếu nhiệm vụ \(i\)
thuộc tập hiện tại. Nhưng nếu định nghĩa như vậy, ta sẽ không viết được
phương trình chuyển trạng thái, vì còn liên quan đến độ dài đường ngắn nhất ở
``điểm nối'' giữa hai nhiệm vụ trước sau. Vì thế ta cần thêm một đại lượng
\(x\) để giới hạn nhiệm vụ đầu tiên cần thực hiện trong tập nhiệm vụ hiện tại.
Một trạng thái được định nghĩa là
\[
  (x,t_1,t_2,\ldots,t_N),
\]
trong đó \(t_x\) bắt buộc bằng 1. Trọng số ứng với trạng thái là độ dài lộ
trình ngắn nhất để hoàn thành tập nhiệm vụ do trạng thái mô tả. Phương trình
chuyển là
\[
  (x,t_1,t_2,\ldots,t_N)
  = \min \{(y,p_1,p_2,\ldots,p_N)+\mathrm{Dis}(x,y)\}.
\]
Trong đó \(y\ne x\) và \(t_y=1\); nếu \(i\ne x\) thì \(p_i=t_i\), còn
\(p_x=0\). Giả sử \(\mathrm{Dis}(x,y)\) là khoảng cách ngắn nhất từ thành phố
kết thúc của nhiệm vụ \(x\) đến thành phố bắt đầu của nhiệm vụ \(y\). Điều
kiện đầu của chuyển trạng thái là
\[
  (x,t_1,t_2,\ldots,t_N)
  = \{\text{khoảng cách từ thành phố kết thúc của nhiệm vụ }x\text{ đến }A\},
\]
trong đó \(t_x=1\), các \(t_i\) còn lại đều bằng 0.

Trong xử lý thực tế, dùng một mảng tối đa 12 chiều như vậy rõ ràng không tiện.
Vì mỗi bit chỉ có thể là 0 hoặc 1, ta xem \((t_1,t_2,\ldots,t_N)\) là biểu
diễn nhị phân của một số thập phân. Khi đó mảng trạng thái là một mảng hai
chiều. Chương trình cụ thể xem Pro\_4\_1.Pas trong phụ lục.

Từ bài toán này có thể thấy, chọn trạng thái trong quy hoạch động thực chất
là chọn cách mô tả thích đáng và ngắn gọn nhất cho sự vật trong bài toán.
Việc chọn trạng thái không thể hoàn thành trong một bước, mà là quá trình dần
điều chỉnh và hoàn thiện trong khi suy nghĩ.

\subsection{Lưu trữ trạng thái}

Sau khi trạng thái đã được chọn, vấn đề tiếp theo là lưu trữ trạng thái. Về
lý thuyết, mỗi trạng thái đều nên lưu hai giá trị: trọng số tối ưu của trạng
thái ấy và giá trị đánh dấu quyết định. Nhưng trên thực tế, nhiều bài quy
hoạch động có số trạng thái rất lớn; muốn chương trình hóa thuật toán, ta
phải tối ưu lưu trữ. Tối ưu chủ yếu là bỏ mọi lượng lưu trữ không cần thiết.

Trong một số bài, đề bài chỉ cho quan hệ trong một phạm vi nào đó, nên ta chỉ
cần lưu các trạng thái trong phạm vi ấy. Xét bài toán sau.

\paragraph{Ví dụ 5.}
Cấu trúc của một sinh vật có thể biểu diễn bằng một dãy ``nguyên tố'', và một
nguyên tố được biểu diễn bằng một chuỗi chữ cái tiếng Anh. Với một tập nguyên
tố \(P\), có thể tạo chuỗi \(S\) bằng cách nối lần lượt \(N\) nguyên tố
\(P_1,P_2,\ldots,P_N\). Bài toán là: cho chuỗi \(S\) và tập nguyên tố \(P\),
tìm tiền tố dài nhất của \(S\) có thể tạo bởi các nguyên tố trong \(P\). Độ
dài lớn nhất của mỗi nguyên tố là 20, độ dài lớn nhất của chuỗi là 500000. Ví
dụ tập nguyên tố là
\(\{A,AB,BBC,CA,BA\}\), chuỗi là \texttt{ABABACABAABCB}, thì độ dài lớn nhất
là 11; cách ghép cụ thể xem Hình 7, trong đó các nguyên tố khác nhau được tô
màu khác nhau.

Trạng thái của bài toán này rất dễ xác định. Mỗi vị trí trong chuỗi chỉ có
hai tính chất: tiền tố của chuỗi kết thúc tại vị trí ấy có thể được tạo bởi
dãy nguyên tố hay không. Vì vậy ta đặt một mảng
\[
  \texttt{Could: Array[1..Max] Of Boolean}
\]
để biểu diễn trạng thái của từng vị trí. Phương trình chuyển là
\[
\begin{split}
  \mathrm{Could}[K] := \mathrm{Could}[K]\ \mathrm{Or}\ (&
  \mathrm{Could}[K-W]\ \mathrm{And}\\
  &S[K-W+1\ldots K]=\mathrm{Ji}[I]).
\end{split}
\]
Trong đó giá trị đầu của \(\mathrm{Could}[k]\) đều là \texttt{False};
\(\mathrm{Ji}[I]\) biểu diễn chuỗi ứng với nguyên tố thứ \(I\), độ dài là
\(W\); \(S[K-W+1\ldots K]\) là chuỗi gồm các ký tự từ vị trí \(K-W+1\) đến
\(K\) trong chuỗi đã cho.

Đến đây, bài toán dường như đã được giải quyết: chỉ cần tìm \(I\) lớn nhất
sao cho \(\mathrm{Could}[I]\) là \texttt{True}. Nhưng quay lại nhìn đề bài, ta
thấy chuỗi dài nhất là 500000, căn bản không thể lưu trữ theo cách thô. Phải
làm gì? Quan sát kỹ phương trình chuyển, ta thấy \(\mathrm{Could}[I]\) chỉ
liên quan đến \(\mathrm{Could}[I-W]\), trong khi \(W\) là độ dài nguyên tố và
độ dài nguyên tố lớn nhất là 20. Nói cách khác, muốn suy ra
\(\mathrm{Could}[I]\), ta nhiều nhất chỉ cần biết các giá trị từ
\(\mathrm{Could}[I-20]\) đến \(\mathrm{Could}[I-1]\). Các giá trị khác có thể
bỏ qua. Như vậy số giá trị \(\mathrm{Could}\) cần ghi giảm đột ngột từ 500000
xuống 20; thay đổi này là rất lớn. Chương trình cụ thể xem Pro\_5.Pas.

Vì sao bài toán này chỉ cần ghi rất ít trạng thái? Đó là do phương trình
chuyển và thứ tự giải đặc biệt của nó. Thứ tự giải được quyết định theo thứ
tự vị trí trong chuỗi, còn trong phương trình chuyển, trạng thái hiện tại chỉ
liên quan đến 20 trạng thái đã giải trước nó. Vì vậy bài toán này có thể xử
lý như trên.

Nhưng với đa số bài toán, cấu trúc đẹp như vậy chưa chắc tồn tại. Khi giải
những bài ấy, ta chỉ có thể cố hết sức giảm lượng lưu trữ. Xét ví dụ sau.

\paragraph{Ví dụ 6.}
Một công ty nhập vào một lô thùng, tổng số là \(N\), \(1\le N\le 1000\), được
đưa vào lần lượt bằng băng chuyền, rồi xếp trong kho thành nhiều nhất \(P\)
cột, \(1\le P\le 4\), như Hình 8.

Biết rằng các thùng nhiều nhất thuộc \(M\) loại, \(1\le M\le 20\). Người ta
muốn các thùng cùng loại được xếp gần nhau nhất có thể để đẹp mắt. Độ đẹp
\(T\) được định nghĩa là:
\[
  T=\sum(\text{số loại khác nhau nhìn thấy lần lượt trong mỗi cột}).
\]
``Số loại khác nhau nhìn thấy lần lượt'' nghĩa là nếu thùng thứ \(K\) trong
một cột khác loại với thùng thứ \(K-1\), thì giá trị độ đẹp tăng thêm 1. Hãy
tìm một cách điều phối sao cho tổng độ đẹp của các cột là nhỏ nhất.

Ý tưởng của bài toán như sau. Thùng thứ \(K\) cần chọn một cột để đặt. Nếu nó
được đặt lên một thùng khác loại, nó làm tổng độ đẹp tăng thêm 1. Nói cách
khác, vị trí đặt thùng thứ \(K\) chỉ liên quan đến tình hình ở đỉnh của \(P\)
cột hiện tại. Nếu muốn tìm giá trị nhỏ nhất sau khi đặt thùng này, tự nhiên
ta cần giá trị nhỏ nhất của \(K-1\) thùng trước đó trong tình hình đỉnh cột
tương ứng. Bài toán quyết định nhiều giai đoạn này đồng thời thỏa nguyên lý
tối ưu và nguyên tắc không hậu hiệu, nên có thể giải bằng quy hoạch động.

Cách chọn trạng thái rất rõ ràng: \((k,P_1,P_2,P_3,P_4)\). \(K\) là số hiệu
của thùng sắp được đặt; \(P_1\) đến \(P_4\) là tình hình đỉnh của mỗi cột sau
khi đặt thùng thứ \(K\), tức loại thùng đang nằm trên đỉnh mỗi cột, giả sử
hiện có bốn cột. Cách chia giai đoạn cũng rõ ràng: chia theo thứ tự thùng đi
tới từ băng chuyền. Phương trình chuyển trạng thái là
\[
  (K,P_1,P_2,P_3,P_4)
  = \min (K-1,Q_1,Q_2,Q_3,Q_4)+\mathrm{Dis}.
\]
Trong đó \(Q_1\) đến \(Q_4\) là tình hình các cột trước khi đặt thùng thứ
\(K\), còn \(\mathrm{Dis}\) bằng 0 hoặc 1, tùy thùng thứ \(K\) có khác loại
với thùng ngay dưới nó trong cột hay không.

Đến đây, mọi bước của thuật toán quy hoạch động đã rõ, dường như có thể viết
chương trình trực tiếp. Nhưng với bài này, lưu trữ một lượng lớn thông tin là
vấn đề bắt buộc phải đối diện. Không gian khả dụng lớn nhất mà Pascal cho phép
là 640KB; trong chế độ bảo vệ, nhiều nhất cũng chỉ hơn 1000KB. Một chương
trình quy hoạch động thường cần lưu hai giá trị: thông tin quyết định và giá
trị trạng thái. Với bài này, nếu làm theo cách thông thường, lúc nhiều nhất
cần \(1000\times 20^3\times (1+2)\) byte, hoàn toàn không thể lưu. Ta phải
tối ưu lưu trữ và bỏ mọi lượng không cần thiết:
\begin{enumerate}
  \item Chỉ ghi thông tin quyết định ở mỗi bước, không ghi giá trị trạng thái
  ở mỗi bước; nhờ vậy có thể tiết kiệm \(\frac{2}{3}\) không gian.
  \item Với trạng thái ở mỗi bước, \(P\) đỉnh cột tương ứng không chứa hai
  thùng cùng loại.
  \item Ngoài cột được chọn ở bước này, thứ tự sắp xếp của \(P-1\) đỉnh cột
  còn lại không ảnh hưởng đến kết quả; vì vậy trạng thái cần lưu là tổ hợp
  của \(P-1\) cột còn lại, không phải hoán vị.
\end{enumerate}

Như vậy tổng lượng lưu trữ có thể giảm từ 24000KB xuống 1200KB, giảm khoảng
20 lần, và hoàn toàn thao tác được trong chế độ bảo vệ của Pascal. Chương
trình cụ thể xem Pro\_6.Pas.

Qua hai bài trên, ta thấy trong ứng dụng thực tế, tốc độ của quy hoạch động
hầu như không thành vấn đề; nhưng khâu lưu trữ trạng thái thì phải suy nghĩ
đi suy nghĩ lại, tối ưu rồi tối ưu nữa. Trong các đề thi vài năm gần đây, khó
điểm và trọng điểm của đa số bài quy hoạch động chính là chọn và lưu trữ
trạng thái. Đây cũng là trọng điểm hàng đầu trong ứng dụng thực tế của quy
hoạch động.

\section{Suy nghĩ sâu hơn về quy hoạch động}

Trên đây đã thảo luận nền tảng lý thuyết của quy hoạch động và các vấn đề có
thể gặp trong ứng dụng thực tế. Vậy quy hoạch động tốt ở đâu? Tại sao ta cần
dùng quy hoạch động? Từ các ví dụ trên có thể thấy ưu điểm lớn nhất của quy
hoạch động là tiết kiệm thời gian. Ưu điểm rất lớn này có thể thấy từ cách
giải Ví dụ 1. Trong thực tế, thời gian của thuật toán tìm kiếm và thuật toán
quy hoạch động khác nhau rất lớn. Bảng 1 là thời gian chạy của hai chương
trình tương ứng với Ví dụ 1.

\[
\begin{array}{c|ccccc}
\text{Dữ liệu thử} & 1 & 2 & 3 & 4 & 5\\ \hline
\text{Tìm kiếm} & 1.04\text{s} & 1.42\text{s} & 1.05\text{s} &
7.42\text{s} & 22.92\text{s}\\
\text{Quy hoạch động} & <0.01\text{s} & <0.01\text{s} & 0.06\text{s} &
0.02\text{s} & 0.01\text{s}
\end{array}
\]
\[
\text{Bảng 1}
\]

Từ bảng thời gian chạy này có thể thấy, thời gian chạy của quy hoạch động và
tìm kiếm trong ứng dụng thực tế khác nhau rất lớn. Về lý thuyết, độ phức tạp
thời gian của tìm kiếm là \(O(N!)\), còn độ phức tạp thời gian của quy hoạch
động là \(O(N^2)\); hai thời gian này căn bản không cùng một cấp. Tăng trưởng
thời gian cấp mũ là rất đáng sợ, lớn hơn rất nhiều so với thuật toán cấp đa
thức. Vì vậy, ta thường gọi thuật toán cấp đa thức là thuật toán ``khá tốt'',
và thuật toán cấp mũ là thuật toán ``khá kém''.

Về bản chất, chương trình quy hoạch động là một kỹ thuật dùng chi phí không
gian để đổi lấy chi phí thời gian. Trong quá trình cài đặt, nó buộc phải lưu
các trạng thái sinh ra. Vì vậy, độ phức tạp không gian của nó lớn hơn các
thuật toán khác. Lấy bài toán hành trình bitonic làm ví dụ: bài này cũng có
thể giải bằng tìm kiếm, xem chương trình Pro\_3\_2.Pas. Độ phức tạp thời gian
của quy hoạch động là \(O(N^2)\), của tìm kiếm là \(O(N!)\); nhưng xét độ
phức tạp không gian, quy hoạch động là \(O(N^2)\), còn tìm kiếm là \(O(N)\).
Tìm kiếm lại tốt hơn quy hoạch động về không gian. Vậy vì sao ta chọn quy
hoạch động để giải Ví dụ 3? Vì không gian của quy hoạch động có thể chịu
được, thực tế là rất dễ chịu, còn thời gian của tìm kiếm thì quá lớn. Vì vậy
ta bỏ không gian để lấy thời gian.

Điều này còn rõ hơn trong Ví dụ 4, tức bài toán vận chuyển bò. Về bản chất,
Ví dụ 4 là bài toán tìm chu trình Hamilton tối ưu. Đây là một bài toán NP; đến
nay người ta vẫn chưa tìm được lời giải cấp đa thức cho nó. Vì vậy, khi dùng
quy hoạch động, ta buộc phải ghi lại từng trạng thái rồi truy hồi. Khi truy
hồi, dường như ta đang thực hiện một thuật toán cấp đa thức. Thực chất, số
trạng thái đã ghi lại vốn đã là cấp mũ. Thuật toán quy hoạch động của ta, về
lý thuyết, có độ phức tạp thời gian gần như thuật toán tìm kiếm thông thường,
còn độ phức tạp không gian lại cao hơn tìm kiếm thông thường rất nhiều: không
gian của quy hoạch động là \(O(2^N)\), của tìm kiếm là \(O(N)\). Vậy vì sao
dùng quy hoạch động cho bài này lại nhanh hơn tìm kiếm? Nhanh hơn bao nhiêu
có thể xem Bảng 2. Phần không gian dùng thêm rốt cuộc để làm gì? Câu trả lời
cho hai câu hỏi này chính là mục đích căn bản của kỹ thuật quy hoạch động:
xử lý dư thừa.

Trong quá trình dùng tìm kiếm, ta làm một số việc đã từng làm. Lấy Ví dụ 4:
giả sử tập nhiệm vụ hiện cần giải là \([1,2,3,4]\). Nếu ta chọn hai nhiệm vụ
đầu tiên cần giải là \(1,2\), thì còn phải chọn thứ tự giải từ hai nhiệm vụ
sau là \(3,4\). Nếu ta chọn hai nhiệm vụ đầu tiên là \(2,1\), khác với
\(1,2\) vì thứ tự giải khác nhau, thì lại phải giải lại công việc ``chọn thứ
tự giải từ hai nhiệm vụ sau \(3,4\)''. Lãng phí lớn biết bao! Câu này đã được
nhắc một lần khi phân tích Ví dụ 1 ở đầu bài; nhắc lại ở đây nhằm vạch rõ tư
duy cơ bản của kỹ thuật quy hoạch động: không làm lại việc đã làm.

Với những việc đã làm, ta đã biết kết quả của chúng. Khi gặp lại cùng một
công việc lần thứ hai, hoàn toàn không cần làm lại; chỉ cần lấy kết quả lần
trước ra dùng. Ý tưởng ``lười biếng'' này chính là thái độ căn bản của con
người khi nghiên cứu vấn đề trong suốt hàng nghìn năm. Các định lý, định luật
trong toán học, vật lý và các ngành tự nhiên mà ta thường gặp, về bản chất,
chẳng phải đều là ``bản ghi công việc'' hay sao? Mục đích học các định lý,
định luật ấy một mặt là học phương pháp suy nghĩ, mặt khác chẳng phải là để
lần sau gặp những vấn đề đó thì giải trực tiếp hay sao? Vì vậy, kỹ thuật quy
hoạch động thực chất là sự tiếp nối của ý tưởng ấy trong lĩnh vực máy tính,
một sự bộc lộ tự nhiên.

Mục đích của việc giải bài là để lần sau không phải giải bài nữa. Nói theo
một câu cổ, điều ta hy vọng cuối cùng là: khi giải bài có thể ``lấy cái bất
biến ứng với vạn biến'', vì vạn biến vẫn không rời khỏi gốc, và bài toán
``gốc'' ấy ta đã từng giải; lần này chỉ cần đọc dữ liệu và đọc kết quả.

Hiểu rõ đạo lý này, năng lực giải bài của ta sẽ tiến thêm một bậc. Với Ví dụ
4, ta hoàn toàn có thể dùng thuật toán tìm kiếm, đồng thời dùng một mảng ghi
lại mọi trạng thái để hỗ trợ tìm kiếm: một khi phát hiện nhiệm vụ tìm kiếm
này đã từng được tìm, ta trực tiếp đọc giá trị tối ưu. Như vậy không còn dư
thừa, thời gian của thuật toán tự nhiên sẽ tốt như quy hoạch động, thậm chí
còn tốt hơn vì trong tìm kiếm có thêm cắt nhánh giữa chừng. Xem Bảng 2.

\[
\begin{array}{c|ccccc}
\text{Dữ liệu thử} & 1 & 2 & 3 & 4 & 5\\ \hline
\text{Quy hoạch động} & 0.11\text{s} & 0.09\text{s} & 0.22\text{s} &
0.22\text{s} & 0.60\text{s}\\
\text{Tìm kiếm} & 1.58\text{s} & 1.90\text{s} & 5.27\text{s} &
6.58\text{s} & 7.09\text{s}\\
\text{Tìm kiếm cải tiến} & 0.06\text{s} & 0.01\text{s} & 0.05\text{s} &
0.07\text{s} & 0.16\text{s}
\end{array}
\]
\[
\text{Bảng 2}
\]

Từ đây ta phát hiện: tìm kiếm chậm vì dư thừa; sau khi bỏ phần dư thừa, công
việc trở nên đơn giản hơn nhiều. Tạm thời không xét yếu tố có thể cắt nhánh
giữa chừng, tức mọi công việc trong quá trình tìm kiếm đều có thể sinh ra kết
quả ta cần. Khi đó những công việc này là bắt buộc, mọi thuật toán đều phải
làm, không thể thiếu, và cũng là giới hạn công việc để giải bài toán tổng
thể. Vì không làm thêm bất kỳ việc nào, tốc độ giải bài của ta sao có thể
không nhanh?

Đó là nguyên nhân căn bản khiến quy hoạch động có tốc độ đáng kinh ngạc: nếu
không đổi hẳn hướng suy nghĩ, khối lượng công việc ta làm sẽ không ít hơn
giới hạn ấy: xét mỗi trạng thái khả dĩ đúng một lần.

\section{Kết luận}

Qua các thảo luận trên, ta biết rằng quy hoạch động là một kỹ thuật có hiệu
suất rất cao. Thời gian chạy nhỏ là vì nó lưu kết quả của những công việc đã
làm. Đa số bài toán quyết định nhiều giai đoạn hầu như đều có thuật toán quy
hoạch động; chỉ có một số bài bắt buộc phải ghi mọi trạng thái, gây căng
thẳng về không gian, nên cuối cùng phải từ bỏ quy hoạch động.

Vấn đề thời gian gần như không thành vấn đề đối với thuật toán quy hoạch
động, vì thuật toán quy hoạch động đã đạt tới giới hạn thời gian. Nhưng vấn đề
không gian lại là vấn đề lớn nhất khi ứng dụng kỹ thuật quy hoạch động. Lượng
lưu trữ trạng thái khổng lồ là điểm đau đầu nhất khi dùng quy hoạch động, và
tự nhiên trở thành một điểm kiểm tra quy hoạch động trong thi tin học. Chọn
trạng thái thế nào, tối ưu lưu trữ trạng thái ra sao, lại càng là vấn đề mỗi
người có cách nhìn riêng. Cuối cùng, tôi tổng kết các bước đại khái để giải
bài quy hoạch động:

\begin{enumerate}
  \item \term{Suy nghĩ} (Think): nghĩ ra phương pháp giải bài. Bước này đặt
  trên cơ sở hiểu và phân tích bài toán; thông qua việc nghiền ngẫm cẩn thận
  bài toán, phân biệt nghiêm túc ưu khuyết của từng hướng, rồi chọn hướng
  hiệu quả và thích đáng nhất để giải bài.
  \item \term{Lập kế hoạch} (Design): xác định các điều kiện của quy hoạch
  động: trạng thái, giai đoạn, quyết định, phương trình chuyển, điều kiện
  biên. Đây là bước điều phối toàn cục, đặt nền tảng thuật toán quy hoạch
  động và điều kiện áp dụng cho bài toán.
  \item \term{Thực hiện} (Operate): chuyển thuật toán thành chương trình. Kết
  hợp với điều kiện máy tính và môi trường lập trình để giải quyết cụ thể vấn
  đề lưu trữ trạng thái, từ đó giải toàn bộ bài toán. Bước này được hoàn
  thành nhờ sự phối hợp giữa con người và máy tính.
\end{enumerate}

Theo các bước trên, một chương trình quy hoạch động về cơ bản đã được hoàn
thành.

Suy cho cùng, quy hoạch động chỉ là một kỹ thuật, cần được hoàn thiện và phát
triển qua ứng dụng liên tục. Chỉ thông qua không ngừng ứng dụng và suy nghĩ,
ta mới có thể hiểu rõ cấu trúc của nó, lĩnh hội nội hàm của nó, và thật sự
nắm vững kỹ thuật này.

\appendix
\section{Phụ lục}

Bài viết có 6 ví dụ. Đề bài trong phần chính đều được viết ngắn; ở đây đưa ra
mô tả chi tiết và chương trình giải cho từng ví dụ. Các hình được nhắc đến
trong đề đã xuất hiện ở phần chính, với cùng số hiệu.

\subsection{Ví dụ 1}

\paragraph{Đề bài.}
Hình 1 cho một mạng đường đi. Số trên đường nối giữa hai điểm biểu diễn
khoảng cách hoặc chi phí giữa hai điểm. Hãy tìm một lộ trình từ \(A\) đến
\(E\) sao cho tổng khoảng cách hoặc tổng chi phí nhỏ nhất.

\paragraph{Ghi chú.}
Bài này lấy từ \textit{Hướng dẫn thi Olympic tin học: thuật toán và thiết kế
chương trình trong tổ hợp}. Trong phần chính, để so sánh hiệu suất của tìm
kiếm và quy hoạch động, bài viết đưa ra hai chương trình Pro\_1\_1.Pas và
Pro\_1\_2.Pas, cùng năm dữ liệu thử Data\_1\_i.Txt và đáp án tương ứng
Sol\_1\_i.Txt trong thư mục Data.

\paragraph{Chương trình.}
\begin{verbatim}
Program Pro_1_1;

Const
  Max        =300;
  Inputfile  ='Input.Txt';
  Outputfile ='Ouptut.Txt';
  Big        =1000000;

Type
  Maps       =Array [1..Max] Of Integer;

Var
  Se         :Set Of Byte;
  Map        :Array[1..Max] Of ^maps;
  Qiu,Pr     :Array[1..Max] Of Byte;
  F          :Text;
  Lp         :Integer;
  N,M        :Integer;
  Min        :Longint;

Procedure Init;
Var
  I,J,K,W    :Integer;
Begin
  Assign(F,Inputfile);
  Reset(F);
  Readln(F,N,M);
  For I:=1 To N Do
  Begin
    New(Map[I]);
    Fillchar(Map[I]^,Sizeof(Map[I]^),0);
  End;
  For I:=1 To M Do
  Begin
    Readln(F,J,K,W);
    Map[J]^[K]:=W;
  End;
  Close(F);
End;

Procedure Search(Dep:Byte;Al:Longint);
Var
  I          :Byte;
Begin
  If Al>=Min Then Exit;
  If Qiu[Dep-1]=N Then
  Begin
    Min:=Al;
    Pr:=Qiu;
    Lp:=Dep-1;
  End
  Else
    For I:=2 To N Do
      If (I In Se) And (Map[Qiu[Dep-1]]^[I]>0) Then
      Begin
        Exclude(Se,I);
        Qiu[Dep]:=I;
        Search(Dep+1,Al+Map[Qiu[Dep-1]]^[Qiu[Dep]]);
        Include(Se,I);
      End;
End;

Procedure Print;
Var
  I          :Integer;
Begin
  Assign(F,Outputfile);
  Rewrite(F);
  Writeln(F,Min);
  For I:=1 To Lp Do
    Write(Pr[I],' ');
  Close(F);
End;

Begin
  Init;
  Se:=[2..N];
  Qiu[1]:=1;
  Min:=Big;
  Search(2,0);
  Print;
End.

Program Pro_1_2;

Const
  Max        =300;
  Inputfile  ='input.Txt';
  Outputfile ='output.Txt';
  Big        =1000000;

Type
  Maps       =Array[1..Max] Of Integer;

Var
  Se         :Set Of Byte;
  Map        :Array[1..Max] Of ^maps;
  Dis        :Array[1..Max] Of Longint;
  Fr         :Array[1..Max] Of Byte;
  Bo         :Array[1..Max] Of Boolean;
  N,M        :Integer;
  F          :Text;

Procedure Init;
Var
  I,J,K,W    :Integer;
Begin
  Assign(F,Inputfile);
  Reset(F);
  Readln(F,N,M);
  For I:=1 To N Do
  Begin
    New(Map[I]);
    Fillchar(Map[I]^,Sizeof(Map[I]^),0);
  End;
  For I:=1 To M Do
  Begin
    Readln(F,J,K,W);
    Map[J]^[K]:=W;
  End;
  Close(F);
End;

Procedure Main;
Var
  I,J,Who :Integer;
  Min     :Longint;
Begin
  Dis[N]:=0;
  Who:=N;
  Fillchar(Fr,Sizeof(Fr),0);
  Fillchar(Bo,Sizeof(Bo),True);
  Bo[N]:=False;
  While Who<>1 Do
  Begin
    For I:=1 To N Do
      If Map[I]^[Who]>0 Then
        If (Fr[I]=0) Or (Dis[I]>Dis[Who]+Map[I]^[Who]) Then
        Begin
          Dis[I]:=Dis[Who]+Map[I]^[Who];
          Fr[I]:=Who;
        End;
    Min:=Big;
    For I:=1 To N Do
      If Bo[I] And (Fr[I]>0) And (Dis[I]<Min) Then
      Begin
        Who:=I;
        Min:=Dis[I];
      End;
    Bo[Who]:=False;
  End;
End;

Procedure Print;
Var
  I        :Integer;
Begin
  Assign(F,Outputfile);
  Rewrite(F);
  Writeln(F,Dis[1]);
  I:=1;
  While I<>N Do
  Begin
    Write(F,I,' ');
    I:=Fr[I];
  End;
  Writeln(F,N);
  Close(F);
End;

Begin
  Init;
  Main;
  Print;
End.
\end{verbatim}

\subsection{Ví dụ 2}

\paragraph{Đề bài.}
Có bốn điểm \(A,B,C,D\), như Hình 4. Giữa hai điểm kề nhau có hai đường nối
\(C_{2k},C_{2k-1}\) \((1\le k\le 3)\), biểu diễn hai đường có thể đi. Số trên
đường nối biểu diễn độ dài đường. Ta định nghĩa trong mọi đường từ \(A\) đến
\(D\), đường có số dư nhỏ nhất khi độ dài chia cho 4 là đường tối ưu.

\paragraph{Ghi chú.}
Bài này lấy từ tạp chí \textit{Olympic tin học}, số 1--2 năm 1998.

\subsection{Ví dụ 3}

\paragraph{Đề bài.}
Bài toán người bán hàng Euclid là bài toán xác định một hành trình khép kín
liên kết \(N\) điểm cho trước trên mặt phẳng. Hình 5(a) cho lời giải của bài
toán bảy điểm. Dạng tổng quát của bài toán này là NP-đầy đủ, do đó lời giải
cần nhiều hơn thời gian đa thức.

J. L. Bentley đề nghị đơn giản hóa bằng cách chỉ xét hành trình bitonic:
xuất phát từ điểm trái nhất, đi nghiêm ngặt từ trái sang phải đến điểm phải
nhất, rồi nghiêm ngặt từ phải sang trái về điểm xuất phát. Hình 5(b) biểu
diễn hành trình bitonic ngắn nhất của cùng bảy điểm. Trong trường hợp này,
một thuật toán thời gian đa thức là có thể.

Hãy mô tả một thuật toán \(O(N^2)\) để xác định hành trình bitonic tối ưu. Có
thể giả sử mọi cặp nút đều có hoành độ khác nhau.

\paragraph{Ghi chú.}
Bài này lấy từ \textit{Đại toàn tập cấu trúc dữ liệu máy tính và thuật toán
thực dụng}. Bài viết trình bày hai cách giải: quy hoạch động và tìm kiếm.

\paragraph{Chương trình.}
\begin{verbatim}
Program Pro_3_1;
Const
  Inputfile  ='input.Txt';
  Outputfile ='Output.Txt';
  Max        =100;
  Big        =10000;

Var
  F          :Text;
  Po         :Array[1..Max,1..2] Of Integer;
  Dis        :Array[1..Max,1..Max] Of Real;
  N          :Integer;

Procedure Init;
Var
  I          :Integer;
Begin
  Assign(F,Inputfile);
  Reset(F);
  Readln(F,N);
  For I:=1 To N Do
    Readln(F,Po[I,1],Po[I,2]);
  Close(F);
End;

Function Len(P1,P2:Integer):Real;
Begin
  Len:=Sqrt(Sqr(Po[P1,1]-Po[P2,1])+Sqr(Po[P1,2]-Po[P2,2]));
End;

Procedure Main;
Var
  I,J,K :Integer;
  Now   :Real;
Begin
  Dis[N,N]:=0;
  For I:=N-1 Downto 1 Do
  Begin
    Dis[I,N]:=Len(I,I+1)+Dis[I+1,N];
    Dis[N,I]:=Dis[I,N];
  End;

  For I:=N-2 Downto 1 Do
    For J:=N-1 Downto I+1 Do
    Begin
      If I+1<J Then Now:=Dis[I+1,J]+Len(I,I+1)
      Else
      Begin
        Now:=Big;
        For K:=J+1 To N Do
          If Dis[K,J]+Len(I,K)<Now Then
            Now:=Dis[K,J]+Len(I,K);
      End;
      For K:=J+1 To N Do
        If Dis[I,K]+Len(J,K)<Now Then
          Now:=Dis[I,K]+Len(J,K);
      Dis[I,J]:=Now;
      Dis[J,I]:=Now;
    End;
End;

Procedure Print;
Var
  X1,X2,I,D :Integer;
  Now       :Real;
  P         :Array[1..2] Of Integer;
  Pr        :Array[1..2,1..Max] Of Byte;
  Ok        :Boolean;

  Procedure Change;
  Var
    G       :Integer;
  Begin
    D:=3-D;
    G:=X1;X1:=X2;X2:=G;
  End;

Begin
  Assign(F,Outputfile);
  Rewrite(F);
  Now:=Big;
  For I:=2 To N Do
    If Dis[1,I]+Len(1,I)<Now Then
    Begin
      Now:=Dis[1,I]+Len(1,I);
      X2:=I;
    End;
  X1:=1;
  Writeln(F,Now:0:2);
  P[1]:=1;P[2]:=1;
  Pr[1,1]:=X1;Pr[2,1]:=X2;
  D:=1;
  While (X1<>N) And (X2<>N) Do
  Begin
    Ok:=True;
    If X1+1<X2 Then
    Begin
      If Dis[X1+1,X2]+Len(X1,X1+1)=Dis[X1,X2] Then
      Begin
        Inc(X1);
        Inc(P[D]);
        Pr[D,P[D]]:=X1;
        Ok:=False;
      End;
    End
    Else
    Begin
      I:=X2+1;
      While (I<=N) And (Dis[I,X2]+Len(X1,I)<>Dis[X1,X2]) Do
        Inc(I);
      If I<=N Then
      Begin
        Ok:=False;
        X1:=I;
        Inc(P[D]);
        Pr[D,P[D]]:=X1;
        Change;
      End;
    End;
    If Ok Then
    Begin
      I:=X2+1;
      While Dis[X1,I]+Len(I,X2)<>Dis[X1,X2] Do
        Inc(I);
      X2:=I;
      Inc(P[3-D]);
      Pr[3-D,P[3-D]]:=X2;
    End;
  End;
  While Pr[D,P[D]]<>N Do
  Begin
    Inc(P[D]);
    Pr[D,P[D]]:=Pr[D,P[D]-1]+1;
  End;

  For I:=1 To P[1] Do
    Write(F,Pr[1,I],' ');
  Writeln(F);
  For I:=1 To P[2] Do
    Write(F,Pr[2,I],' ');
  Close(F);
End;

Begin
  Init;
  Main;
  Print;
End.

Program Pro_3_2;
Const
  Inputfile  ='input.Txt';
  Outputfile ='Output.Txt';
  Max        =100;
  Big        =10000;

Var
  F          :Text;
  Po         :Array[1..Max,1..2] Of Integer;
  Qiu,Pr     :Array[1..Max] Of Byte;
  Lpr        :Integer;
  N          :Integer;
  Min        :Real;

Procedure Init;
Var
  I          :Integer;
Begin
  Assign(F,Inputfile);
  Reset(F);
  Readln(F,N);
  For I:=1 To N Do
    Readln(F,Po[I,1],Po[I,2]);
  Close(F);
End;

Function Len(P1,P2:Integer):Real;
Begin
  Len:=Sqrt(Sqr(Po[P1,1]-Po[P2,1])+Sqr(Po[P1,2]-Po[P2,2]));
End;

Procedure Search(Dep,Fr,Last:Byte;Al:Real);
Var
  I,P :Byte;
  Rs  :Real;

  Procedure Did;
  Var
    K :Byte;
  Begin
    K:=Fr+1;
    Rs:=0;
    P:=Last;
    While K<>I Do
    Begin
      Rs:=Rs+Len(P,K);
      P:=K;
      Inc(K);
    End;
    Rs:=Rs+Len(Fr,I);
  End;

Begin
  If Al>=Min Then Exit;
  If Fr=N Then
  Begin
    If Al+Len(Last,N)<Min Then
    Begin
      Min:=Al+Len(Last,N);
      Pr:=Qiu;
      Lpr:=Dep-1;
    End;
  End
  Else
    For I:=Fr+1 To N Do
    Begin
      Qiu[Dep]:=I;
      Did;
      Search(Dep+1,I,P,Al+Rs);
    End
End;

Procedure Print;
Var
  I          :Byte;
  Se         :Set Of Byte;
Begin
  Assign(F,Outputfile);
  Rewrite(F);
  Writeln(F,Min:0:2);
  Se:=[1..N];
  For I:=1 To Lpr Do
  Begin
    Write(F,Pr[I],' ');
    Se:=Se-[Pr[I]];
  End;
  Se:=Se+[1,N];
  Writeln(F);
  For I:=1 To N Do
    If I In Se Then
      Write(F,I,' ');
  Close(F);
End;

Begin
  Init;
  Min:=Big;
  Search(1,1,1,0);
  Print;
End.
\end{verbatim}

\subsection{Ví dụ 4}

\paragraph{Đề bài.}
Có một công ty vận chuyển bò đặt tại nông trại \(A\), có một xe dùng để vận
chuyển bò giữa bảy nông trại \(A,B,C,D,E,F,G\). Khoảng cách giữa bảy nông
trại được cho trong bảng sau:
\[
\begin{array}{c|rrrrrr}
 & B & C & D & E & F & G\\ \hline
A & 56 & 43 & 71 & 35 & 41 & 36\\
B & 43 & 54 & 58 & 36 & 79 & 31\\
C & 71 & 58 & 30 & 20 & 31 & 58\\
D & 35 & 36 & 20 & 38 & 59 & 75\\
E & 41 & 79 & 31 & 59 & 44 & 67\\
F & 36 & 31 & 58 & 75 & 67 & 72
\end{array}
\]

Mỗi buổi sáng, công ty phải quyết định thứ tự vận chuyển bò sao cho tổng
quãng đường nhỏ nhất. Quy tắc:
\begin{enumerate}
  \item Xe luôn xuất phát từ trụ sở công ty ở nông trại \(A\), và sau khi mọi
  nhiệm vụ hoàn thành phải quay về trụ sở.
  \item Xe mỗi lần chỉ chở một con bò.
  \item Mỗi nhiệm vụ vận chuyển được cho bởi một cặp chữ cái, lần lượt biểu
  diễn bò sẽ được chở từ nông trại nào đến nông trại nào.
\end{enumerate}
Nhiệm vụ của bạn là viết chương trình, với một tập nhiệm vụ cho trước, tìm
lộ trình ngắn nhất để hoàn thành mọi nhiệm vụ, xuất phát từ nông trại \(A\)
và cuối cùng trở về nông trại \(A\).

\paragraph{Dữ liệu vào.}
Các nhiệm vụ đặt trong tệp \texttt{Delvr.Txt}. Dòng đầu là tổng số nhiệm vụ
\(N\), \(1\le N\le 12\). Từ dòng thứ hai, mỗi nhiệm vụ gồm một cặp chữ cái
cách nhau bởi một dấu cách; chữ cái đầu biểu diễn nông trại nơi con bò đang
ở, chữ cái thứ hai biểu diễn nông trại đích. Chương trình phải chạy không quá
60 giây.

\paragraph{Ví dụ vào.}
\begin{verbatim}
5
FC
GB
BD
AE
GA
\end{verbatim}

\paragraph{Dữ liệu ra.}
In ra màn hình độ dài lộ trình ngắn nhất, và hiển thị thứ tự các nông trại mà
xe đi qua.

\paragraph{Ví dụ ra.}
\begin{verbatim}
368
AEFCGBDGA
\end{verbatim}

\paragraph{Ghi chú.}
Bài này lấy từ tạp chí \textit{Olympic tin học}, số 1--2 năm 1998. Trong phần
chính, bài này có ba thuật toán: quy hoạch động, tìm kiếm, và tìm kiếm cải
tiến; lần lượt ứng với các chương trình Pro\_4\_1.Pas, Pro\_4\_2.Pas,
Pro\_4\_3.Pas. Ngoài ra có năm dữ liệu thử; tệp vào và ra là
Data\_4\_i.Txt, Sol\_4\_i.Txt, đặt trong thư mục Data.

\paragraph{Chương trình.}
\begin{verbatim}
Program Pro_4_1;
Const
  Map:Array['A'..'G','A'..'G'] Of Byte=
    ((0,56,43,71,35,41,36),
     (56,0,54,58,36,79,31),
     (43,54,0,30,20,31,58),
     (71,58,30,0,38,59,75),
     (35,36,20,38,0,44,67),
     (41,79,31,59,44,0,72),
     (36,31,58,75,67,72,0));

  Inputfile  ='d1.Inp';
  Outputfile ='sol_4_1.Txt';

Type
  Kus         =Array[0..4095] Of Word;
  Dirs        =Array[0..4095] Of Byte;
  Ses         =Set Of 1..12;

Var
  F           :Text;
  N           :Integer;
  Wu          :Array[0..12,1..2] Of Char;
  Dis         :Array[0..12,0..12] Of Integer;
  Ku          :Array[1..12] Of ^kus;
  Dir         :Array[1..12] Of ^dirs;

Procedure Init;
Var
  I           :Integer;
  Ch          :Char;
Begin
  Assign(F,Inputfile);
  Reset(F);
  Readln(F,N);
  For I:=1 To N Do
    Readln(F,Wu[I,1],Ch,Wu[I,2]);
  Close(F);
End;

Procedure Prepare;
Var
  I,J         :Integer;
Begin
  Wu[0,1]:='A';Wu[0,2]:='A';
  For I:=0 To N Do
    For J:=0 To N Do
      Dis[I,J]:=Map[Wu[I,2],Wu[J,1]];
End;

Procedure Main;
Var
  Last,I,K,What :Word;
  S             :Ses;

  Function Num(S:Ses):Word;
  Var
    X       :Word Absolute S;
  Begin
    Num:=X Div 2;
  End;

  Procedure Did(Dep,From:Byte;S:Ses);
  Var
    I       :Byte;
  Begin
    If Dep>K Then
    Begin
      For I:=1 To N Do
        If (I In S) And
           (Ku[I]^[Num(S-[I])]+Dis[What,I]<Ku[What]^[Num(S)]) Then
        Begin
          Ku[What]^[Num(S)]:=Ku[I]^[Num(S-[I])]+Dis[What,I];
          Dir[What]^[Num(S)]:=I;
        End;
    End
    Else
      For I:=From+1 To N Do
        If I<>What Then
          Did(Dep+1,I,S+[I]);
  End;

Begin
  For I:=1 To N Do
  Begin
    New(Ku[I]);
    New(Dir[I]);
    Fillchar(Ku[I]^,Sizeof(Ku[I]^),$ff);
  End;
  For I:=1 To N Do
    Ku[I]^[0]:=Dis[I,0];

  For K:=1 To N-1 Do
    For What:=1 To N Do
      Did(1,0,[]);

  S:=[1..N];
  K:=60000;
  For I:=1 To N Do
    If Dis[0,I]+Ku[I]^[Num(S-[I])]<K Then
    Begin
      K:=Dis[0,I]+Ku[I]^[Num(S-[I])];
      What:=I;
    End;
  For I:=1 To N Do
    Inc(K,Map[Wu[I,1],Wu[I,2]]);

  Assign(F,Outputfile);
  Rewrite(F);
  Writeln(F,K);
  Write(F,'A ');
  Last:=0;
  For I:=1 To N Do
  Begin
    If Wu[Last,2]<>Wu[What,1] Then Write(F,Wu[What,1],' ');
    Write(F,Wu[What,2],' ');
    Exclude(S,What);
    Last:=What;
    What:=Dir[What]^[Num(S)];
  End;
  If 'A'<>Wu[Last,2] Then Write(F,'A');
  Close(F);
End;

Begin
  Init;
  Prepare;
  Main;
End.

Program Pro_4_2;

Const
  Map:Array['A'..'G','A'..'G'] Of Byte=
    ((0,56,43,71,35,41,36),
     (56,0,54,58,36,79,31),
     (43,54,0,30,20,31,58),
     (71,58,30,0,38,59,75),
     (35,36,20,38,0,44,67),
     (41,79,31,59,44,0,72),
     (36,31,58,75,67,72,0));
  Inputfile  ='d1.Inp';
  Outputfile ='data_4_1.Txt';

Var
  F           :Text;
  N           :Integer;
  Min         :Word;
  Wu          :Array[0..12,1..2] Of Char;
  Dis         :Array[0..12,0..12] Of Integer;
  Qiu,Pr      :Array[0..14] Of Byte;
  Se          :Set Of Byte;

Procedure Init;
Var
  I           :Integer;
  Ch          :Char;
Begin
  Assign(F,Inputfile);
  Reset(F);
  Readln(F,N);
  For I:=1 To N Do
    Readln(F,Wu[I,1],Ch,Wu[I,2]);
  Close(F);
End;

Procedure Prepare;
Var
  I,J         :Integer;
Begin
  Wu[0,1]:='A';Wu[0,2]:='A';
  For I:=0 To N Do
    For J:=0 To N Do
      Dis[I,J]:=Map[Wu[I,2],Wu[J,1]];
End;

Procedure Search(Dep:Byte;Al:Word);
Var
  I           :Byte;
Begin
  If Al>=Min Then Exit;
  If Dep=N+1 Then
  Begin
    If Dis[Qiu[N],0]+Al<Min Then
    Begin
      Min:=Dis[Qiu[N],0]+Al;
      Pr:=Qiu;
    End;
  End
  Else
    For I:=1 To N Do
      If I In Se Then
      Begin
        Qiu[Dep]:=I;
        Exclude(Se,I);
        Search(Dep+1,Al+Dis[Qiu[Dep-1],I]);
        Include(Se,I);
      End;
End;

Procedure Print;
Var
  I           :Byte;
Begin
  For I:=1 To N Do
    Inc(Min,Map[Wu[I,1],Wu[I,2]]);
  Assign(F,Outputfile);
  Rewrite(F);
  Writeln(F,Min);
  Write(F,'A ');
  For I:=1 To N Do
  Begin
    If Wu[Pr[I-1],2]<>Wu[Pr[I],1] Then Write(F,Wu[Pr[I],1],' ');
    Write(F,Wu[Pr[I],2],' ');
  End;
  If 'A'<>Wu[Pr[N],2] Then Write(F,'A');
  Close(F);
End;

Begin
  Init;
  Prepare;
  Se:=[1..12];
  Qiu[0]:=0;
  Min:=60000;
  Search(1,0);
  Print;
End.

Program Pro_4_3;

Const
  Map:Array['A'..'G','A'..'G'] Of Byte=
    ((0,56,43,71,35,41,36),
     (56,0,54,58,36,79,31),
     (43,54,0,30,20,31,58),
     (71,58,30,0,38,59,75),
     (35,36,20,38,0,44,67),
     (41,79,31,59,44,0,72),
     (36,31,58,75,67,72,0));
  Inputfile  ='d1.Inp';
  Outputfile ='data_4_1.Txt';

Type
  Kus         =Array[0..4095] Of Word;
  Ses         =Set Of 1..12;

Var
  F           :Text;
  N           :Integer;
  Min         :Word;
  Big,I       :Word;
  Wu          :Array[0..12,1..2] Of Char;
  Dis         :Array[0..12,0..12] Of Integer;
  Ku          :Array[1..12] Of ^kus;
  Qiu         :Array[1..13] Of Byte;
  Se          :Ses;

Procedure Init;
Var
  I           :Integer;
  Ch          :Char;
Begin
  Assign(F,Inputfile);
  Reset(F);
  Readln(F,N);
  For I:=1 To N Do
    Readln(F,Wu[I,1],Ch,Wu[I,2]);
  Close(F);
End;

Procedure Prepare;
Var
  I,J         :Integer;
Begin
  Wu[0,1]:='A';Wu[0,2]:='A';
  For I:=0 To N Do
    For J:=0 To N Do
      Dis[I,J]:=Map[Wu[I,2],Wu[J,1]];

  For I:=1 To N Do
  Begin
    New(Ku[I]);
    Fillchar(Ku[I]^,Sizeof(Ku[I]^),$ff);
  End;
  Big:=Ku[1]^[1];
End;

Function Num(S:Ses):Word;
Var
  X           :Word Absolute S;
Begin
  Num:=X Div 2;
End;

Procedure Search(Dep:Byte;Al:Word);
Var
  I           :Byte;
  D           :Word;
Begin
  If Al>=Min Then Exit;
  If Ku[Qiu[Dep-1]]^[Num(Se)]<Big Then
  Begin
    If Al+Ku[Qiu[Dep-1]]^[Num(Se)]<Min Then
      Min:=Al+Ku[Qiu[Dep-1]]^[Num(Se)];
  End
  Else
    If Dep>N Then
    Begin
      Ku[Qiu[Dep-1]]^[Num(Se)]:=Dis[Qiu[Dep-1],0];
      If Al+Ku[Qiu[Dep-1]]^[Num(Se)]<Min Then
        Min:=Al+Ku[Qiu[Dep-1]]^[Num(Se)];
    End
    Else
    Begin
      D:=Big;
      For I:=1 To N Do
        If I In Se Then
        Begin
          Exclude(Se,I);
          Qiu[Dep]:=I;
          Search(Dep+1,Al+Dis[Qiu[Dep-1],Qiu[Dep]]);
          If Ku[I]^[Num(Se)]+Dis[Qiu[Dep-1],Qiu[Dep]]<D Then
            D:=Ku[I]^[Num(Se)]+Dis[Qiu[Dep-1],Qiu[Dep]];
          Include(Se,I);
        End;
      Ku[Qiu[Dep-1]]^[Num(Se)]:=D;
    End;
End;

Procedure Print;
Var
  I,J,Last    :Integer;
  D           :Word;
Begin
  Se:=[1..N];
  J:=1;
  While Dis[0,J]+Ku[J]^[Num(Se-[J])]<>Min Do
    Inc(J);
  D:=Min;
  For I:=1 To N Do
    Inc(D,Map[Wu[I,1],Wu[I,2]]);
  Assign(F,Outputfile);
  Rewrite(F);
  Writeln(F,D);
  Write(F,'A ');
  Dec(Min,Dis[0,J]);
  Last:=0;

  For I:=1 To N Do
  Begin
    If Wu[J,1]<>Wu[Last,2] Then
      Write(F,Wu[J,1],' ');
    Write(F,Wu[J,2],' ');
    If I<>N Then
    Begin
      Last:=J;
      Exclude(Se,J);
      J:=1;
      While (Not (J In Se)) Or
            (Ku[J]^[Num(Se-[J])]+Dis[Last,J]<>Min) Do
        Inc(J);
      Dec(Min,Dis[Last,J]);
    End;
  End;
  If Wu[J,2]<>'A' Then Write(F,'A');
  Close(F);
End;

Begin
  Init;
  Prepare;
  Se:=[1..N];
  Min:=Big;
  For I:=1 To N Do
  Begin
    Se:=[1..N]-[I];
    Qiu[1]:=I;
    Search(2,Dis[0,I]);
  End;
  Print;
End.
\end{verbatim}

\subsection{Ví dụ 5}

\paragraph{Đề bài.}
Một số cấu trúc phức tạp của sinh vật có thể biểu diễn bằng dãy các nguyên
tố, còn một nguyên tố được biểu diễn bằng một chuỗi chữ cái in hoa. Một vấn
đề của các nhà sinh học là phân giải một dãy dài như vậy thành một dãy các
nguyên tố. Với tập nguyên tố \(P\) cho trước, nếu có thể chọn \(N\) nguyên tố
\(P_1,P_2,P_3,\ldots,P_n\) từ đó, nối các chuỗi tương ứng theo thứ tự để được
một chuỗi \(S\), thì gọi \(S\) có thể được tạo bởi tập nguyên tố \(P\). Khi
chọn nguyên tố từ \(P\), một nguyên tố có thể dùng nhiều lần hoặc không dùng.
Ví dụ, dãy \texttt{aBaBaCaBaaB} có thể tạo bởi tập nguyên tố
\(\{\texttt{a},\texttt{aB},\texttt{Ba},\texttt{Ca},\texttt{BBC}\}\).

\(K\) ký tự đầu của một chuỗi được gọi là tiền tố của chuỗi đó, với độ dài
\(K\). Hãy viết chương trình: với tập nguyên tố \(P\) và chuỗi \(T\) nhập
vào, tìm một tiền tố của \(T\) có thể tạo bởi tập nguyên tố \(P\), sao cho độ
dài tiền tố lớn nhất có thể, và xuất độ dài ấy.

\paragraph{Dữ liệu vào.}
Có hai tệp vào \texttt{INPUT.TXT} và \texttt{DATA.TXT}. Dòng đầu của
\texttt{INPUT.TXT} là số nguyên tố \(N\) trong tập \(P\),
\(1\le N\le 100\). Sau đó có \(2N\) dòng, mỗi hai dòng mô tả một nguyên tố:
dòng đầu là độ dài \(L\), \(1\le L\le 20\); dòng sau là chuỗi chữ cái in hoa
độ dài \(L\), biểu diễn nguyên tố ấy. Mọi nguyên tố đôi một khác nhau.

Tệp \texttt{DATA.TXT} mô tả chuỗi \(T\). Mỗi dòng bắt đầu bằng một chữ cái in
hoa; dòng cuối là ký tự \texttt{.}, biểu thị kết thúc chuỗi. Độ dài của \(T\)
ít nhất là 1 và nhiều nhất không vượt quá 500000.

\paragraph{Dữ liệu ra.}
Tệp ra tên \texttt{OUTPUT.TXT}. Chỉ có một dòng là một số, biểu thị độ dài
tiền tố dài nhất của \(T\) có thể tạo bởi \(P\).

\paragraph{Ghi chú.}
Bài này là đề thi Olympic Tin học Quốc tế lần thứ tám. Chương trình ở dưới.
Dữ liệu vào và ra mẫu là Data\_5.Txt, Input\_5.Txt, Sol\_5.Txt trong thư mục
Data.

\paragraph{Chương trình.}
\begin{verbatim}
Program Pro_5;

Const
  Datafile   ='Data.Txt';
  Inputfile  ='Input.Txt';
  Outputfile ='Output.Txt';

Type
  Jis = Record
          Lst : Byte;
          St  : String[20];
        End;

Var
  F          :Text;
  Ji         :Array[1..100] Of Jis;
  Ch         :Array[0..21] Of Char;
  Bo         :Array[0..21] Of Boolean;
  N          :Integer;
  Tot        :Longint;
  Big        :Longint;

Procedure Init;
Var
  I          :Integer;
Begin
  Assign(F,Inputfile);
  Reset(F);
  Readln(F,N);
  For I:=1 To N Do
  Begin
    Readln(F,Ji[I].Lst);
    Readln(F,Ji[I].St);
  End;
  Close(F);
End;

Procedure Main;
Var
  I,J        :Integer;

  Procedure See;
  Var
    I        :Byte;
    X,Y      :Integer;
  Begin
    Bo[J]:=False;
    For I:=1 To N Do
    Begin
      X:=J;Y:=Ji[I].Lst;
      While (Y>0) And (Ji[I].St[Y]=Ch[X]) Do
      Begin
        Dec(X);
        Dec(Y);
        If X<0 Then X:=21;
      End;
      If (Y=0) And (Bo[X]) Then
      Begin
        Bo[J]:=True;
        Exit;
      End;
    End;
  End;

Begin
  Assign(F,Datafile);
  Reset(F);
  Big:=0;
  Tot:=0;
  Bo[0]:=True;
  J:=0;
  While (Tot-Big<20) And (Ch[J]<>'.') Do
  Begin
    Inc(Tot);Inc(J);
    If J>21 Then J:=0;
    Readln(F,Ch[J]);
    If Ch[J]<>'.' Then
    Begin
      See;
      If Bo[J] Then Big:=Tot;
    End;
  End;
  Close(F);
End;

Procedure Print;
Begin
  Assign(F,Outputfile);
  Rewrite(F);
  Writeln(F,Big);
  Close(F);
End;

Begin
  Init;
  Main;
  Print;
End.
\end{verbatim}

\subsection{Ví dụ 6}

\paragraph{Đề bài.}
Một công ty nhập vào một lô thùng, tổng số là \(N\), \(1\le N\le 1000\), được
đưa vào lần lượt bằng băng chuyền, rồi được xếp trong kho thành nhiều nhất
\(P\) cột, \(1\le P\le 4\). Quá trình xếp này do một máy điều phối thực hiện,
như Hình 8.

Biết các thùng nhập vào nhiều nhất thuộc \(M\) loại, \(1\le M\le 20\). Ông
chủ tự nhiên muốn các thùng được xếp càng đẹp càng tốt, tức là xếp các thùng
cùng loại gần nhau nhất có thể. Theo sở thích của ông chủ, độ đẹp \(T\) được
định nghĩa là
\[
  T=\sum(\text{số loại khác nhau nhìn thấy lần lượt trong mỗi cột}).
\]
``Số loại khác nhau nhìn thấy lần lượt'' nghĩa là nếu thùng thứ \(K\) khác
loại với thùng thứ \(K-1\), thì tính thêm một loại. Ví dụ, cột
\texttt{1 2 3 1} có số loại khác nhau nhìn thấy lần lượt là 4; cột
\texttt{1 2 2 3 3 1 2} có số loại này là 5. Vì vậy \(T\) càng nhỏ thì càng
đẹp, và ông chủ càng hài lòng.

Đã biết không hạn chế số thùng trong mỗi cột; có thể xếp mọi thùng vào một
cột, cũng có thể xếp thành \(P\) cột. Nhưng do thùng được vận chuyển bằng
băng chuyền nên thứ tự thùng là cố định: nếu thùng \(A\) đến trước thùng
\(B\), thì nếu \(A,B\) được xếp cùng một cột, \(A\) phải nằm trước \(B\).

Hãy viết một chương trình cho máy điều phối, để máy dựa vào thông tin thùng
do băng chuyền gửi tới mà điều phối, sao cho sau khi xếp thành \(P\) cột thì
độ đẹp nhỏ nhất.

\paragraph{Dữ liệu vào.}
Tệp vào tên \texttt{INPUT.TXT}. Nội dung gồm hai dòng. Dòng đầu có ba số theo
thứ tự là \(N,M,P\). Dòng thứ hai có \(N\) số, là loại của các thùng do băng
chuyền đưa tới; số thứ \(K\) là loại của thùng thứ \(K\).

\paragraph{Dữ liệu ra.}
Tệp ra tên \texttt{OUTPUT.TXT}. Nội dung có hai dòng. Dòng đầu là độ đẹp
\(T\). Dòng thứ hai có \(N\) số, là một phương án điều phối; số thứ \(K\) biểu
thị thùng thứ \(K\) sẽ đi đến cột nào.

\paragraph{Ghi chú.}
Bài này do tác giả tự đặt. Chương trình ở dưới; dữ liệu vào và ra mẫu là
Data\_6.Txt, Sol\_6.Txt trong thư mục Data.

\paragraph{Chương trình.}
\begin{verbatim}
Program Pro_6;
Const
  Fin = 'Input.txt';
  Fou = 'Output.txt';

Type
  Kus = Array[0..1160] Of Byte;
  Zis = Array[0..1160] Of Integer;

Var
  Ku       : Array[1..1000] Of ^kus;
  Shu,Shu1 : Array[1..1000] Of Byte;
  Zi       : Array[0..1] Of Zis;
  Use1     : Array[0..20,0..20,0..20] Of Integer;
  Use2     : Array[0..1160,1..3] Of Byte;
  Total,I,M,N,P,N1 : Integer;
  F        : Text;

Procedure Init;
Var
  I,J      :Integer;
Begin
  Assign(F,Fin);
  Reset(F);
  Readln(F,N1,M,P);
  For I:=1 To N1 Do
    Read(F,Shu1[I]);
  Shu[1]:=Shu1[1];
  N:=1;
  For I:=2 To N1 Do
    If Shu1[I]<>Shu1[I-1] Then
    Begin
      Inc(N);
      Shu[N]:=Shu1[I];
    End;
  Close(F);
End;

Procedure Diduse;
Var
  Qs       :Array[1..3] Of Byte;
  Stop     :Byte;
  I,J,K    :Byte;

  Procedure Serch(Depth,From:Byte);
  Var
    I      :Byte;
  Begin
    If Depth=Stop+1 Then
    Begin
      Inc(Total);
      Use2[Total,1]:=Qs[1];Use2[Total,2]:=Qs[2];
      Use2[Total,3]:=Qs[3];
      Use1[Qs[1],Qs[2],Qs[3]]:=Total;
      Use1[Qs[1],Qs[3],Qs[2]]:=Total;
      Use1[Qs[2],Qs[1],Qs[3]]:=Total;
      Use1[Qs[2],Qs[3],Qs[1]]:=Total;
      Use1[Qs[3],Qs[1],Qs[2]]:=Total;
      Use1[Qs[3],Qs[2],Qs[1]]:=Total;
    End
    Else
      For I:=From To M-1 Do
      Begin
        Qs[Depth]:=I;
        Serch(Depth+1,I+1);
      End;
  End;

Begin
  Total:=0;
  Fillchar(Qs,Sizeof(Qs),0);
  For I:=0 To M Do
    For J:=0 To M Do
      For K:=0 To M Do
        Use1[I,J,K]:=-1;
  Use2[0,1]:=0;Use2[0,2]:=0;Use2[0,3]:=0;
  Use1[0,0,0]:=0;
  For Stop:=1 To P-1 Do
    Serch(1,1);
End;

Procedure Main;
Var
  I,J,K,W1,W2,W3 :Integer;
  Q               :Array[1..4] Of Byte;
Begin
  For I:=1 To N Do
  Begin
    New(Ku[I]);
  End;
  Fillchar(Zi,Sizeof(Zi),0);
  Fillchar(Ku[1]^,Sizeof(Ku[1]^),0);
  Zi[0][0]:=1;
  Fillchar(Q,Sizeof(Q),0);
  For I:=2 To N Do
  Begin
    For J:=0 To Total Do
      If Zi[0][J]<>0 Then
      Begin
        W1:=0;
        For K:=1 To 3 Do
        Begin
          Q[K]:=Use2[J,K];
          If Q[K]>=Shu[I-1] Then Inc(Q[K]);
          If Q[K]=Shu[I] Then W1:=K;
        End;
        Q[4]:=Shu[I-1];
        If W1<>0 Then
        Begin
          W2:=Q[4];Q[4]:=Q[W1];Q[W1]:=W2;
          For K:=1 To 3 Do
            If Q[K]>Shu[I] Then Dec(Q[K]);
          W3:=Use1[Q[1],Q[2],Q[3]];
          If (Zi[1][W3]=0) Or (Zi[1][W3]>Zi[0][J]) Then
          Begin
            Zi[1][W3]:=Zi[0][J];
            Ku[I]^[W3]:=Shu[I];
          End;
        End
        Else
        Begin
          For K:=1 To 4 Do
            If Q[K]>Shu[I] Then Dec(Q[K]);
          For K:=1 To 4 Do
          Begin
            Case K Of
              1:W1:=Use1[Q[2],Q[3],Q[4]];
              2:W1:=Use1[Q[1],Q[3],Q[4]];
              3:W1:=Use1[Q[1],Q[2],Q[4]];
              4:W1:=Use1[Q[1],Q[2],Q[3]];
            End;
            If W1<>-1 Then
              If (Zi[1][W1]=0) Or (Zi[1][W1]>Zi[0][J]+1) Then
              Begin
                Zi[1][W1]:=Zi[0][J]+1;
                W2:=Q[K];
                If W2>=Shu[I] Then Inc(W2);
                Ku[I]^[W1]:=W2;
              End;
          End;
        End;
      End;
    Zi[0]:=Zi[1];
    Fillchar(Zi[1],Sizeof(Zi[1]),0);
  End;
End;

Procedure Print;
Var
  Pr,Pr1          :Array[1..1000] Of Integer;
  Now,Q           :Array[1..4] Of Byte;
  I,J,K,Ps,W1,W2  :Integer;
Begin
  Ps:=1001;
  For J:=0 To Total Do
    If (Zi[0][J]<>0) And (Zi[0][J]<Ps) Then
    Begin
      Ps:=Zi[0][J];
      K:=J;
    End;
  Assign(F,Fou);
  Rewrite(F);
  Writeln(F,Ps);
  Now[2]:=Use2[K,1];
  Now[3]:=Use2[K,2];
  Now[4]:=Use2[K,3];
  For I:=2 To 4 Do
    If Now[I]>=Shu[N] Then Inc(Now[I]);
  Now[1]:=Shu[N];
  Pr[N]:=1;
  For K:=N-1 Downto 1 Do
  Begin
    J:=1;
    While Now[J]<>Shu[K+1] Do
      Inc(J);
    W2:=0;
    For I:=1 To 4 Do
      If I<>J Then
      Begin
        Inc(W2);
        Q[W2]:=Now[I];
        If Q[W2]>=Shu[K+1] Then Dec(Q[W2]);
      End;
    W1:=Use1[Q[1],Q[2],Q[3]];
    Now[J]:=Ku[K+1]^[W1];
    For I:=1 To 4 Do
      If Now[I]=Shu[K] Then
        Pr[K]:=I;
  End;
  J:=1;
  Pr1[1]:=Pr[1];
  For I:=2 To N1 Do
  Begin
    If Shu1[I]<>Shu1[I-1] Then Inc(J);
    Pr1[I]:=Pr[J];
  End;
  For I:=1 To N1 Do
    Write(F,Pr1[I],' ');
  Close(F);
End;

Begin
  Init;
  Diduse;
  Main;
  Print;
End.
\end{verbatim}

Các chương trình trên đều được kiểm thử thành công trong môi trường Pascal
7.0.

\begin{thebibliography}{9}
\bibitem{comb}
\textit{Hướng dẫn thi Olympic tin học quốc tế và toàn quốc dành cho thanh
thiếu niên: thuật toán và thiết kế chương trình trong tổ hợp}, Wu Wenhu và
Wang Jiande biên soạn, Nhà xuất bản Đại học Thanh Hoa.

\bibitem{journal}
Tạp chí \textit{Olympic tin học}, số 1--2 năm 1998.

\bibitem{algorithms}
\textit{Đại toàn tập cấu trúc dữ liệu máy tính và thuật toán thực dụng}, Mou
Ren chủ biên, Công ty máy tính Hy vọng Bắc Kinh.

\bibitem{ioi1996}
\textit{Phân tích đề thi Olympic tin học trong nước và quốc tế năm 1996}, Wu
Wenhu và Wang Jiande biên soạn, Nhà xuất bản Đại học Bắc Kinh.
\end{thebibliography}

\end{document}
